% ============================================================
% 附录
% ============================================================
\ctexset{appendix/name={附录,}, appendix/number={\Alph{chapter}}}
\renewcommand{\chaptermark}[1]{\markboth{附录\,\thechapter\quad #1}{}}

\chapter{课件勘误汇总}
编写本讲义时逐页审读了选择性必修1的七份课件。课件整体质量较好，以下是发现的笔误、不规范表述或与现行教材不一致之处（正文中均在对应位置作了说明，多数以“课件勘误”栏目呈现）：

\begin{center}
\small
\begin{tabularx}{\linewidth}{p{2.6cm}XXp{1.2cm}}
\toprule
课件位置 & 原表述 & 更正 & 所在节 \\
\midrule
1.1 反应热，第47页例1 & \ce{N2(g) + 3H2(aq) = 2NH3(g)} & \ce{H2} 的状态应为 g：\ce{N2(g) + 3H2(g) = 2NH3(g)} & 1.3 节 \\
1.1 反应热，第56页 & 燃烧的指定产物中列有“Cl$\to$\ce{Cl2(g)}” & 教材只规定 C、H、S、N 四种元素的指定产物；含氯物质燃烧时氯一般转化为 HCl，不宜记忆为 \ce{Cl2(g)} & 1.4 节 \\
1.1 反应热，第15页 & “中和热” & 新教材称“中和反应反应热”（含义相同），表述时应与教材一致 & 1.2 节 \\
1.2 反应热的计算，第19页 & “$\Delta H=\sum$反应物的燃烧热$-\sum$生成物的燃烧热” & 只有把燃烧热写成带负号的 $\Delta_\text{c}H$ 时才成立；用正值燃烧热 $Q$ 代入应为 $\Delta H=\sum Q_\text{生}-\sum Q_\text{反}$ & 1.6 节 \\
2.1 化学反应速率，第4页 & “纯液体表示速率可用 L/s 或 L/min” & 规范说法：纯液体、纯固体不能用浓度变化表示速率，可用单位时间内物质的量或质量的变化表示 & 2.1 节 \\
2.2 化学平衡，第38～39页 & 方程式写作“\ce{$m$A(g) + $n$B(g) <=> $e$C(g) + $f$d(g)}” & 应为 \ce{$f$D(g)}（大小写笔误） & 3.2 节 \\
2.3 化学反应方向，第16页 & \ce{Fe(OH)2(s) + O2(g) + H2O(g) = Fe(OH)3(s)}，并称“制备氢氧化亚铁” & 未配平，且是氢氧化亚铁被氧化：\ce{4Fe(OH)2(s) + O2(g) + 2H2O(l) = 4Fe(OH)3(s)} & 3.5 节 \\
\bottomrule
\end{tabularx}
\end{center}

\chapter{常用数据（25 ℃）}
\begin{center}
\small
\begin{tabular}{llll}
\toprule
弱电解质 & 电离常数 & 难溶电解质 & 溶度积 $K_\text{sp}$ \\
\midrule
\ce{CH3COOH} & $K_\text{a}=1.75\times10^{-5}$ & AgCl & $1.8\times10^{-10}$ \\
\ce{NH3*H2O} & $K_\text{b}=1.8\times10^{-5}$ & AgI & $8.5\times10^{-17}$ \\
\ce{H2CO3} & $K_{\text{a}1}=4.5\times10^{-7}$，$K_{\text{a}2}=4.7\times10^{-11}$ & \ce{Ag2CrO4} & $1.1\times10^{-12}$ \\
\ce{H2SO3} & $K_{\text{a}1}=1.4\times10^{-2}$，$K_{\text{a}2}=6.0\times10^{-8}$ & \ce{BaSO4} & $1.1\times10^{-10}$ \\
HClO & $K_\text{a}=4.0\times10^{-8}$ & \ce{BaCO3} & $2.6\times10^{-9}$ \\
HF & p$K_\text{a}=3.2$（2024年广东卷所给） & \ce{CaCO3} & $3.4\times10^{-9}$ \\
HCN & $K_\text{a}=6.2\times10^{-10}$ & \ce{CaSO4} & $4.9\times10^{-5}$ \\
& & \ce{Mg(OH)2} & $5.6\times10^{-12}$ \\
& & \ce{Fe(OH)2} & $4.9\times10^{-17}$ \\
& & \ce{Fe(OH)3} & $2.8\times10^{-39}$ \\
& & \ce{Al(OH)3} & $1.3\times10^{-33}$ \\
& & \ce{Cu(OH)2} & $2.2\times10^{-20}$ \\
\bottomrule
\end{tabular}
\end{center}
\noindent\small 说明：电离常数、溶度积数值随资料来源略有差异，考试时一律以题目所给数据为准。其他常数：$N_\text{A}=6.02\times10^{23}$ mol$^{-1}$，$F=96500$ C·mol$^{-1}$，$R=8.314$ J·mol$^{-1}$·K$^{-1}$，$K_\text{w}=1.0\times10^{-14}$（25 ℃），标准状况下 $V_\text{m}=22.4$ L·mol$^{-1}$。\normalsize

\chapter{高考化学答题规范用语}
\begin{center}
\small
\begin{tabularx}{\linewidth}{p{3.2cm}X}
\toprule
设问 & 规范表述 \\
\midrule
条件选择的原因 & 速率角度：……反应速率较快；平衡角度：……平衡转化率较高；实际角度：催化剂活性最大 / 设备要求 / 成本 / 安全环保 \\
平衡移动解释 & 改变……（条件），使……（平衡）向……方向移动，导致……（结果） \\
沸点或溶解性比较 & ××分子间（与水分子间）能形成氢键 / ××相对分子质量大，范德华力强 / 形成分子内氢键，沸点降低 \\
酸性强弱比较 & ××的电负性大，吸电子能力强，使羧基中 O—H 键极性增强，更易电离出 \ce{H+} \\
滴定终点 & 当滴入最后半滴××溶液时，溶液由×色变为×色，且半分钟内不恢复原色 \\
检验沉淀是否洗净 & 取最后一次洗涤液少许于试管中，滴加××，若无××现象，则已洗净 \\
离子检验 & 取少量待测液于试管中，滴加××，若出现××，则含有××（必要时先排除干扰） \\
提高浸取率 & 粉碎原料、搅拌、适当升温、适当增大浸取液浓度、延长浸取时间 \\
调节 pH 的原因 & 使××完全沉淀而××不沉淀（附 $K_\text{sp}$ 计算） \\
控制温度不宜过高 & 防止××分解（挥发）；防止催化剂失活；平衡逆向移动（放热反应） \\
冰水 / 乙醇洗涤 & 减少晶体溶解损失；乙醇易挥发，便于干燥 \\
\bottomrule
\end{tabularx}
\end{center}
